\documentclass[11pt]{article}

\usepackage[utf8]{inputenc}
\usepackage[T1]{fontenc}
\usepackage[spanish,es-tabla]{babel}
\usepackage{lmodern}
\usepackage[a4paper,margin=2cm]{geometry}
\usepackage{array}
\usepackage{booktabs}
\usepackage{xcolor}
\usepackage{enumitem}
\usepackage{hyperref}

\setlength{\parskip}{0.5em}
\setlength{\parindent}{0pt}
\renewcommand{\arraystretch}{1.3}
\setlist{nosep,leftmargin=1.6em}
\definecolor{marca}{HTML}{0B3A5C}
\definecolor{barra}{HTML}{E6ECF1}
\hypersetup{colorlinks=true,linkcolor=marca,urlcolor=marca}

\newcommand{\urlapp}{https://thankstochatgpt.github.io/LOARD}
\newcommand{\app}[1]{\href{\urlapp#1}{\texttt{#1}}}
\newcommand{\etiqueta}[1]{\colorbox{barra}{\textbf{\color{marca}#1}}}

\title{Historias de usuario del MVP\\[0.3em]
\large Loard - Hito 1}
\author{}
\date{}

\begin{document}
\maketitle
\tableofcontents


Este documento recoge las 12 historias de usuario que componen el MVP (Release 1) del sistema de información para Loard Construcciones S.A.S. En total suman 47 puntos de esfuerzo.

Cada historia indica la pantalla que la resuelve en el prototipo de alta fidelidad del Hito 2, publicado en \href{\urlapp/}{\nolinkurl{https://thankstochatgpt.github.io/LOARD/}}. La correspondencia completa entre historias y pantallas está en la \emph{Cobertura de HU} del Módulo I del Hito 2.

El MVP ataca dos problemas concretos que hoy le generan a la empresa dolores de cabeza y pérdida de tiempo: entregar a tiempo el dossier de calidad que exige Chevron, y evitar que alguien sin la habilitación de SST al día termine asignado a un frente de obra.

\section{Por qué estas historias}

La planeación anterior del proyecto no fue buena y eso nos dejó sin una base clara sobre la cual construir. Para poder avanzar, replanteamos el alcance y definimos estas 12 historias como un MVP acotado. El objetivo es poder avanzar, construir los prototipos y llevarlos a pruebas de usabilidad.

Este no es el backlog definitivo. A medida que iteremos y recojamos lo que salga de las pruebas, iremos agregando más historias de usuario.

\section{Resumen}

\begin{center}\small
\begin{tabular}{>{\raggedright\arraybackslash}p{1.5cm} >{\raggedright\arraybackslash}p{6.1cm} >{\raggedright\arraybackslash}p{4cm} >{\centering\arraybackslash}p{1.7cm} >{\centering\arraybackslash}p{1.2cm}}
\toprule
\textbf{Código} & \textbf{Historia} & \textbf{Rol principal} & \textbf{Prioridad} & \textbf{Puntos} \\ \midrule
HU-M01 & Iniciar sesión y ver solo lo que corresponde al rol & Todos los usuarios & MUST & 3 \\
HU-M02 & Registrar el corte de obra & Ingeniero Residente & MUST & 5 \\
HU-M03 & Adjuntar fotografías y soportes al corte & Ingeniero Residente & MUST & 2 \\
HU-M04 & Visar o devolver el corte & Director de Proyectos & MUST & 3 \\
HU-M05 & Publicar documentos vigentes y armar la checklist & Coordinación SIG / Calidad & MUST & 5 \\
HU-M06 & Cargar soportes en la checklist y buscar documentos & Residente / Calidad & MUST & 3 \\
HU-M07 & Registrar ensayos NDT ligados a una junta & Calidad & MUST & 5 \\
HU-M08 & Monitorear el avance del dossier y exportarlo & Calidad / Gerencia & MUST & 5 \\
HU-M09 & Registrar cursos y aptitud médica con semáforo & Coordinación SST & MUST & 5 \\
HU-M10 & Asignar personal al corte con bloqueo SST & Ingeniero Residente & MUST & 3 \\
HU-M11 & Consultar la bitácora de cambios & Coordinación SIG / Calidad & MUST & 3 \\
HU-M12 & Consultar el tablero de avance & Director / Gerencia & MUST & 5 \\
 & \textbf{Total} &  &  & \textbf{47} \\
\bottomrule
\end{tabular}
\end{center}


\newpage
\section*{HU-M01: Iniciar sesión y ver solo lo que corresponde al rol}
\addcontentsline{toc}{section}{HU-M01: Iniciar sesión y ver solo lo que corresponde al rol}
\etiqueta{MUST}\quad\etiqueta{3 puntos}\quad Pantalla en el prototipo: \app{/acceso}

\begin{tabular}{>{\bfseries}p{2.4cm} p{13.4cm}}
\toprule
Módulo & Control de acceso por roles. \\
\bottomrule
\end{tabular}

\subsection*{Descripción de la historia de usuario}

\begin{tabular}{>{\bfseries}p{2.4cm} p{13.4cm}}
\toprule
Como & Usuario de Loard (Ingeniero Residente, Director de Proyectos, Coordinación SIG/Calidad, Responsable SST o Gerencia). \\
Quiero & iniciar sesión y encontrar únicamente las pantallas que tienen que ver con mi trabajo. \\
Para & no perderme entre opciones que no uso y que nadie toque información que no le corresponde. \\
\bottomrule
\end{tabular}

\subsection*{Criterios de aceptación}

\begin{enumerate}
\item Con usuario y contraseña correctos, el sistema me deja entrar y me lleva a la pantalla de inicio de mi rol.
\item Si los datos están mal, veo un mensaje claro de error que no dice cuál de los dos campos falló.
\item El menú solo muestra los módulos autorizados para mi rol.
\item Si intento abrir por dirección directa una pantalla que no es de mi rol, el sistema no me deja y me avisa que no tengo permiso.
\item No existen cuentas para clientes externos: solo entra personal de Loard.
\item Puedo cerrar sesión cuando quiera y, después de hacerlo, nadie puede volver a las pantallas sin autenticarse de nuevo.
\end{enumerate}


\newpage
\section*{HU-M02: Registrar el corte de obra}
\addcontentsline{toc}{section}{HU-M02: Registrar el corte de obra}
\etiqueta{MUST}\quad\etiqueta{5 puntos}\quad Pantalla en el prototipo: \app{/cortes}

\begin{tabular}{>{\bfseries}p{2.4cm} p{13.4cm}}
\toprule
Módulo & Captura de cortes de obra. \\
\bottomrule
\end{tabular}

\subsection*{Descripción de la historia de usuario}

\begin{tabular}{>{\bfseries}p{2.4cm} p{13.4cm}}
\toprule
Como & Ingeniero Residente. \\
Quiero & registrar desde mi portátil el corte semanal o mensual con el avance físico por actividad, las horas-hombre trabajadas y los insumos consumidos. \\
Para & dejar el avance del frente en un solo lugar y dejar de reportarlo por WhatsApp. \\
\bottomrule
\end{tabular}

\subsection*{Criterios de aceptación}

\begin{enumerate}
\item Puedo crear un corte eligiendo el proyecto, el frente de obra y el periodo que cubre (semanal o mensual).
\item Por cada actividad registro el avance físico, las horas-hombre y los insumos consumidos.
\item El sistema no acepta avances negativos ni acumulados por encima del~100\,\% de la actividad.
\item Puedo guardar el corte como borrador y terminarlo después sin perder lo que llevaba.
\item Cuando lo envío, el corte queda pendiente de visado y le aparece al Director de Proyectos.
\item Tener un corte pendiente de visado no me impide registrar el siguiente.
\end{enumerate}


\newpage
\section*{HU-M03: Adjuntar fotografías y soportes al corte}
\addcontentsline{toc}{section}{HU-M03: Adjuntar fotografías y soportes al corte}
\etiqueta{MUST}\quad\etiqueta{2 puntos}\quad Pantalla en el prototipo: \app{/cortes/ejemplo-3}

\begin{tabular}{>{\bfseries}p{2.4cm} p{13.4cm}}
\toprule
Módulo & Captura de cortes de obra. \\
\bottomrule
\end{tabular}

\subsection*{Descripción de la historia de usuario}

\begin{tabular}{>{\bfseries}p{2.4cm} p{13.4cm}}
\toprule
Como & Ingeniero Residente. \\
Quiero & adjuntar fotografías y soportes en PDF al corte que registré. \\
Para & que lo reportado tenga evidencia y quien lo revise no tenga que pedírmela aparte. \\
\bottomrule
\end{tabular}

\subsection*{Criterios de aceptación}

\begin{enumerate}
\item Puedo adjuntar una o varias fotografías y archivos PDF a un corte.
\item Si subo un archivo de otro tipo o demasiado pesado, el sistema lo rechaza y me explica por qué.
\item Los adjuntos se ven listados dentro del corte y se pueden abrir o descargar.
\item Puedo quitar un adjunto mientras el corte no haya sido visado.
\item El Director de Proyectos ve los mismos adjuntos cuando revisa el corte.
\end{enumerate}


\newpage
\section*{HU-M04: Visar o devolver el corte}
\addcontentsline{toc}{section}{HU-M04: Visar o devolver el corte}
\etiqueta{MUST}\quad\etiqueta{3 puntos}\quad Pantalla en el prototipo: \app{/cortes}

\begin{tabular}{>{\bfseries}p{2.4cm} p{13.4cm}}
\toprule
Módulo & Captura de cortes de obra y visado. \\
\bottomrule
\end{tabular}

\subsection*{Descripción de la historia de usuario}

\begin{tabular}{>{\bfseries}p{2.4cm} p{13.4cm}}
\toprule
Como & Director de Proyectos. \\
Quiero & revisar cada corte y aprobarlo (visarlo) o devolverlo con observaciones puntuales. \\
Para & asegurar que lo reportado es correcto sin frenar el trabajo en el frente. \\
\bottomrule
\end{tabular}

\subsection*{Criterios de aceptación}

\begin{enumerate}
\item Veo la lista de cortes pendientes de visado, con su proyecto, frente y periodo.
\item Al abrir un corte veo el avance, las horas-hombre, los insumos y los adjuntos.
\item Puedo visar el corte, y desde ese momento queda aprobado y ya no se puede modificar.
\item Puedo devolverlo, pero solo si escribo la observación que explica qué hay que corregir.
\item El residente ve el corte devuelto junto con mi observación y puede corregirlo y reenviarlo.
\item Mientras un corte está pendiente o devuelto, el residente puede seguir registrando cortes nuevos en ese frente.
\end{enumerate}


\newpage
\section*{HU-M05: Publicar documentos vigentes y armar la checklist}
\addcontentsline{toc}{section}{HU-M05: Publicar documentos vigentes y armar la checklist}
\etiqueta{MUST}\quad\etiqueta{5 puntos}\quad Pantalla en el prototipo: \app{/dossier}

\begin{tabular}{>{\bfseries}p{2.4cm} p{13.4cm}}
\toprule
Módulo & Dossier de calidad (checklist). \\
\bottomrule
\end{tabular}

\subsection*{Descripción de la historia de usuario}

\begin{tabular}{>{\bfseries}p{2.4cm} p{13.4cm}}
\toprule
Como & Coordinación SIG / Calidad. \\
Quiero & publicar las versiones vigentes de los formatos del SIG y definir los ítems que exige la checklist en cada una de las 6 carpetas contractuales. \\
Para & que todos trabajen con el formato correcto y el dossier se vaya armando desde el primer día, no al cierre del contrato. \\
\bottomrule
\end{tabular}

\subsection*{Criterios de aceptación}

\begin{enumerate}
\item Puedo publicar un documento del SIG indicando su código, nombre y versión.
\item Al publicar una versión nueva, la anterior deja de aparecer como vigente pero se conserva en el historial.
\item Cada proyecto activo tiene su checklist organizada en las 6 carpetas: documentación contractual y legal, actas de corte de obra, planos de ingeniería, certificados de materiales, equipos y NDT, informes técnicos de obra y memorias de cálculo.
\item Puedo agregar, editar o retirar ítems requeridos dentro de cada carpeta.
\item Todo ítem nuevo arranca en estado Pendiente.
\item Solo Coordinación SIG / Calidad puede publicar documentos y modificar la estructura de la checklist.
\end{enumerate}


\newpage
\section*{HU-M06: Cargar soportes en la checklist y buscar documentos}
\addcontentsline{toc}{section}{HU-M06: Cargar soportes en la checklist y buscar documentos}
\etiqueta{MUST}\quad\etiqueta{3 puntos}\quad Pantalla en el prototipo: \app{/dossier/documentos}

\begin{tabular}{>{\bfseries}p{2.4cm} p{13.4cm}}
\toprule
Módulo & Dossier de calidad (checklist). \\
\bottomrule
\end{tabular}

\subsection*{Descripción de la historia de usuario}

\begin{tabular}{>{\bfseries}p{2.4cm} p{13.4cm}}
\toprule
Como & Ingeniero Residente o integrante de Calidad. \\
Quiero & cargar el soporte PDF en cada ítem de la checklist y buscar con filtros los documentos vigentes. \\
Para & ir completando el dossier a medida que avanza la obra y encontrar rápido el formato que necesito. \\
\bottomrule
\end{tabular}

\subsection*{Criterios de aceptación}

\begin{enumerate}
\item Puedo cargar un PDF en un ítem de la checklist y el ítem pasa a estado Archivo Cargado.
\item Calidad puede marcar el ítem como Validado, o como Observado dejando escrito el motivo.
\item Si un ítem queda Observado, puedo reemplazar el archivo y el ítem vuelve a Archivo Cargado.
\item Puedo buscar documentos vigentes filtrando al menos por nombre o código, carpeta y estado.
\item La búsqueda solo me muestra versiones vigentes, nunca las obsoletas.
\item Si no hay resultados, el sistema lo dice con claridad en lugar de mostrar una lista vacía sin explicación.
\end{enumerate}


\newpage
\section*{HU-M07: Registrar ensayos NDT ligados a una junta}
\addcontentsline{toc}{section}{HU-M07: Registrar ensayos NDT ligados a una junta}
\etiqueta{MUST}\quad\etiqueta{5 puntos}\quad Pantalla en el prototipo: \app{/dossier/ensayos}

\begin{tabular}{>{\bfseries}p{2.4cm} p{13.4cm}}
\toprule
Módulo & Dossier de calidad (checklist). \\
\bottomrule
\end{tabular}

\subsection*{Descripción de la historia de usuario}

\begin{tabular}{>{\bfseries}p{2.4cm} p{13.4cm}}
\toprule
Como & Integrante de Calidad. \\
Quiero & registrar cada ensayo no destructivo de laboratorio asociado a la junta soldada que se inspeccionó y vincularlo a la checklist. \\
Para & tener la trazabilidad de cada junta y que el reporte quede de una vez dentro del dossier. \\
\bottomrule
\end{tabular}

\subsection*{Criterios de aceptación}

\begin{enumerate}
\item Puedo registrar un ensayo indicando la junta soldada, el tipo de ensayo, la fecha, el laboratorio y el resultado.
\item El sistema no deja guardar un ensayo sin junta asociada.
\item Puedo adjuntar el reporte PDF que entrega el laboratorio.
\item El ensayo queda vinculado al ítem correspondiente de la carpeta 4 (certificados de materiales, equipos y NDT).
\item Desde una junta puedo ver todos los ensayos que se le han hecho.
\item Un ensayo con resultado rechazado queda visible como tal y no cuenta como requisito cumplido.
\end{enumerate}


\newpage
\section*{HU-M08: Monitorear el avance del dossier y exportarlo}
\addcontentsline{toc}{section}{HU-M08: Monitorear el avance del dossier y exportarlo}
\etiqueta{MUST}\quad\etiqueta{5 puntos}\quad Pantalla en el prototipo: \app{/dossier/avance}

\begin{tabular}{>{\bfseries}p{2.4cm} p{13.4cm}}
\toprule
Módulo & Dossier de calidad (checklist). \\
\bottomrule
\end{tabular}

\subsection*{Descripción de la historia de usuario}

\begin{tabular}{>{\bfseries}p{2.4cm} p{13.4cm}}
\toprule
Como & Integrante de Calidad o de Gerencia. \\
Quiero & ver el porcentaje de avance de las 6 carpetas, saber qué requisitos faltan, controlar los días que quedan del plazo y exportar el paquete validado. \\
Para & radicar el dossier dentro de los 10 días calendario posteriores a la aceptación del acta de facturación. \\
\bottomrule
\end{tabular}

\subsection*{Criterios de aceptación}

\begin{enumerate}
\item Veo el porcentaje de avance de cada carpeta y el global, y se actualizan apenas cambia el estado de un ítem.
\item El sistema resalta los ítems que faltan (Pendiente u Observado) para que sepa exactamente qué perseguir.
\item Al registrar la fecha de aceptación del acta de facturación, arranca la cuenta regresiva de 10 días calendario.
\item Los días restantes se muestran siempre a la vista y la alerta se hace más notoria a medida que se acerca el vencimiento.
\item El cierre y la exportación solo se habilitan cuando el~100\,\% de los ítems está Validado.
\item La exportación genera un archivo .ZIP organizado en las 6 carpetas contractuales.
\end{enumerate}


\newpage
\section*{HU-M09: Registrar cursos y aptitud médica con semáforo}
\addcontentsline{toc}{section}{HU-M09: Registrar cursos y aptitud médica con semáforo}
\etiqueta{MUST}\quad\etiqueta{5 puntos}\quad Pantalla en el prototipo: \app{/sst}

\begin{tabular}{>{\bfseries}p{2.4cm} p{13.4cm}}
\toprule
Módulo & Seguridad y Salud en el Trabajo (SST). \\
\bottomrule
\end{tabular}

\subsection*{Descripción de la historia de usuario}

\begin{tabular}{>{\bfseries}p{2.4cm} p{13.4cm}}
\toprule
Como & Responsable de SST. \\
Quiero & registrar la vigencia de los cursos críticos de cada trabajador y su concepto médico ocupacional, con un semáforo que me avise antes de que algo se venza. \\
Para & renovar a tiempo y que nadie llegue a obra con la habilitación vencida. \\
\bottomrule
\end{tabular}

\subsection*{Criterios de aceptación}

\begin{enumerate}
\item Por cada trabajador registro sus cursos críticos (trabajo en alturas, espacios confinados, inducciones) con fecha de vencimiento.
\item Registro el concepto médico únicamente como Habilitado o No habilitado.
\item El sistema no pide ni guarda historias clínicas ni diagnósticos, por reserva legal.
\item El semáforo muestra verde cuando todo está vigente, una alerta preventiva cuando faltan 30 días o menos, una alerta más fuerte cuando faltan 15 días o menos y rojo cuando ya venció.
\item Un trabajador con algún curso vencido o con concepto No habilitado aparece como no habilitado para obra.
\item Puedo filtrar la lista de personal por estado del semáforo para ver primero a quienes están por vencerse.
\end{enumerate}


\newpage
\section*{HU-M10: Asignar personal al corte con bloqueo SST}
\addcontentsline{toc}{section}{HU-M10: Asignar personal al corte con bloqueo SST}
\etiqueta{MUST}\quad\etiqueta{3 puntos}\quad Pantalla en el prototipo: \app{/cortes/nuevo/personal}

\begin{tabular}{>{\bfseries}p{2.4cm} p{13.4cm}}
\toprule
Módulo & Captura de cortes de obra / SST. \\
\bottomrule
\end{tabular}

\subsection*{Descripción de la historia de usuario}

\begin{tabular}{>{\bfseries}p{2.4cm} p{13.4cm}}
\toprule
Como & Ingeniero Residente. \\
Quiero & asignar al corte de obra el personal que trabajó en el frente, y que el sistema no me deje escoger a quien no esté habilitado. \\
Para & no exponer a la empresa ni al trabajador por asignar a alguien con la habilitación vencida. \\
\bottomrule
\end{tabular}

\subsection*{Criterios de aceptación}

\begin{enumerate}
\item Al asignar personal, veo a cada trabajador con el estado de su semáforo SST.
\item No puedo seleccionar a un trabajador vencido o No habilitado: el bloqueo es terminante y no hay forma de saltárselo.
\item El sistema me dice por qué está bloqueado (curso vencido o concepto médico), sin mostrar información clínica.
\item Sí puedo asignar a quienes están en alerta de 30 o 15 días, y la advertencia queda a la vista.
\item Si un trabajador ya asignado se vence antes de que envíe el corte, el sistema me avisa y no me deja enviarlo con esa persona.
\end{enumerate}


\newpage
\section*{HU-M11: Consultar la bitácora de cambios}
\addcontentsline{toc}{section}{HU-M11: Consultar la bitácora de cambios}
\etiqueta{MUST}\quad\etiqueta{3 puntos}\quad Pantalla en el prototipo: \app{/dossier/bitacora}

\begin{tabular}{>{\bfseries}p{2.4cm} p{13.4cm}}
\toprule
Módulo & Trazabilidad y seguimiento directivo. \\
\bottomrule
\end{tabular}

\subsection*{Descripción de la historia de usuario}

\begin{tabular}{>{\bfseries}p{2.4cm} p{13.4cm}}
\toprule
Como & Coordinación SIG / Calidad. \\
Quiero & consultar el historial de cambios de estado y aprobaciones de un registro. \\
Para & saber quién hizo qué y cuándo, y poder responder con respaldo ante una auditoría. \\
\bottomrule
\end{tabular}

\subsection*{Criterios de aceptación}

\begin{enumerate}
\item Cada cambio de estado o de documento queda guardado con fecha, hora y autor.
\item Desde un corte, un ítem de la checklist o un documento puedo abrir su historial en orden cronológico.
\item El historial muestra el estado anterior, el nuevo y la observación, cuando la hubo.
\item Nadie puede editar ni borrar entradas de la bitácora.
\item Puedo filtrar el historial por rango de fechas y por autor.
\end{enumerate}


\newpage
\section*{HU-M12: Consultar el tablero de avance}
\addcontentsline{toc}{section}{HU-M12: Consultar el tablero de avance}
\etiqueta{MUST}\quad\etiqueta{5 puntos}\quad Pantalla en el prototipo: \app{/dashboard}

\begin{tabular}{>{\bfseries}p{2.4cm} p{13.4cm}}
\toprule
Módulo & Trazabilidad y seguimiento directivo. \\
\bottomrule
\end{tabular}

\subsection*{Descripción de la historia de usuario}

\begin{tabular}{>{\bfseries}p{2.4cm} p{13.4cm}}
\toprule
Como & Director de Proyectos o integrante de Gerencia. \\
Quiero & ver en un tablero el avance físico de cada proyecto frente al tiempo contractual transcurrido, junto con el progreso del dossier. \\
Para & detectar a tiempo los proyectos atrasados y decidir dónde intervenir sin tener que pedir informes. \\
\bottomrule
\end{tabular}

\subsection*{Criterios de aceptación}

\begin{enumerate}
\item El tablero lista los proyectos activos con su porcentaje de avance físico y el porcentaje de tiempo contractual transcurrido.
\item El avance físico se calcula con los cortes visados.
\item Se distingue a simple vista cuándo un proyecto va por debajo de lo que debería según el tiempo transcurrido.
\item Cada proyecto muestra el porcentaje de avance de la checklist del dossier.
\item Puedo entrar a un proyecto para ver el detalle de sus cortes y de sus carpetas.
\item El tablero es solo de consulta: desde ahí no se modifica ningún dato.
\end{enumerate}

% =====================================================================
% Historias nuevas a partir de la validación RITE del Hito 2.
% \historia{código}{título}{prioridad}{puntos}{módulo}{como}{quiero}{para}{criterios}
\newcommand{\historia}[9]{%
  \newpage
  \section*{#1: #2}
  \addcontentsline{toc}{section}{#1: #2}
  \etiqueta{#3}\quad\etiqueta{#4 puntos}\quad Origen: validación RITE del Hito 2

  \begin{tabular}{>{\bfseries}p{2.4cm} p{13.4cm}}
  \toprule
  Módulo & #5 \\
  \bottomrule
  \end{tabular}

  \subsection*{Descripción de la historia de usuario}

  \begin{tabular}{>{\bfseries}p{2.4cm} p{13.4cm}}
  \toprule
  Como & #6 \\
  Quiero & #7 \\
  Para & #8 \\
  \bottomrule
  \end{tabular}

  \subsection*{Criterios de aceptación}

  \begin{enumerate}
  #9
  \end{enumerate}}

\newpage
\section{Historias surgidas de la validación RITE}

Durante la validación RITE del prototipo (Hito 2, 16 de noviembre de 2026), dos residentes de obra y un director de proyectos pidieron funciones que el MVP no cubre. No son fallas del prototipo: son necesidades del trabajo diario. Esta sección las convierte en historias de usuario para el siguiente release (Release 2), como se anunció al inicio de este documento.

Se identifican con el prefijo \textbf{HU-R}. La prioridad sigue el método MoSCoW (MUST, SHOULD, COULD, WON'T) y los puntos son una estimación inicial que se debe ajustar con el equipo. La fuente es la sección \emph{Pedidos para próximas versiones} del \emph{Reporte de validación RITE} (Módulo III del Hito 2).

\begin{center}\small
\begin{tabular}{>{\raggedright\arraybackslash}p{1.5cm} >{\raggedright\arraybackslash}p{6.1cm} >{\raggedright\arraybackslash}p{4cm} >{\centering\arraybackslash}p{1.7cm} >{\centering\arraybackslash}p{1.2cm}}
\toprule
\textbf{Código} & \textbf{Historia} & \textbf{Rol principal} & \textbf{Prioridad} & \textbf{Puntos} \\ \midrule
HU-R01 & Detallar las actividades de tubería y obra civil & Ingeniero Residente & SHOULD & 3 \\
HU-R02 & Registrar la calificación vigente de los soldadores & Coordinación SST & SHOULD & 3 \\
HU-R03 & Registrar los lotes de materiales que llegan al proyecto & Ingeniero Residente & COULD & 3 \\
HU-R04 & Controlar herramientas, equipos y su mantenimiento & Residente / Director & SHOULD & 5 \\
HU-R05 & Hacer la revisión diaria de equipos y herramientas & Director de Proyectos & COULD & 3 \\
HU-R06 & Anotar el clima del día en el corte & Ingeniero Residente & COULD & 1 \\
HU-R07 & Descargar el corte mensual para el cliente & Director / Gerencia & SHOULD & 3 \\
HU-R08 & Ver solo los cortes y documentos de mi obra & Ingeniero Residente & SHOULD & 3 \\
HU-R09 & Registrar un proyecto nuevo y avisar al director & Ingeniero Residente & SHOULD & 2 \\
HU-R10 & Recibir avisos por correo y recordatorios de facturación & Residente / Calidad & SHOULD & 3 \\
HU-R11 & Que el encargado SST de cada obra actualice sus cursos & Encargado SST de obra & COULD & 2 \\
 & \textbf{Total} &  &  & \textbf{31} \\
\bottomrule
\end{tabular}
\end{center}

\subsection*{Pedidos que ya cubre el MVP}

\begin{itemize}
  \item \textbf{Elegir la obra en el tablero} (director): HU-M12 ya permite escoger el proyecto del que se ven los datos.
  \item \textbf{Medir el avance comparando los días cotizados con los transcurridos} (director): HU-M12 compara el avance físico con el tiempo contractual transcurrido.
\end{itemize}

\subsection*{Pedidos fuera de alcance por ahora (WON'T)}

Quedan registrados, pero salen del objetivo del sistema (dossier de calidad y control SST) y no se convierten en historias en este release:

\begin{itemize}
  \item Pasar las horas trabajadas a la nómina.
  \item Llevar la caja y el balance de cada obra.
  \item Módulo de cotizaciones con precios unitarios, revisión del director en línea y exportación a PDF o Excel.
  \item Ver en el tablero el dinero ejecutado y lo facturado.
  \item Conectar el plan de trabajo de MS Project.
\end{itemize}

\historia{HU-R01}{Detallar las actividades de tubería y obra civil}{SHOULD}{3}
{Captura de cortes de obra.}
{Ingeniero Residente.}
{registrar por separado los metros de tubería tendidos, la cantidad de soldaduras y el diámetro de la tubería, y anotar los pasos previos de la obra civil.}
{que el avance refleje el rendimiento real de cada trabajo, porque no rinde igual una tubería de 4 pulgadas que una de 8.}
{\item En las actividades de tubería registro por separado los metros tendidos y la cantidad de soldaduras.
\item Cada actividad de tubería indica su diámetro (4, 6 u 8 pulgadas).
\item En obra civil puedo registrar los pasos previos: excavación, formaleta, relleno, compactación y solado.
\item Cada paso tiene su propia unidad de medida y su cantidad.
\item Los cortes anteriores siguen mostrando sus actividades como se registraron.}

\historia{HU-R02}{Registrar la calificación vigente de los soldadores}{SHOULD}{3}
{Control de habilitación SST.}
{Coordinación SST.}
{registrar la calificación en soldadura de cada soldador, con su vigencia, además de sus cursos.}
{que ningún soldador sin calificación vigente quede asignado a un frente.}
{\item Puedo registrar la calificación de un soldador con su proceso, fecha de expedición, vencimiento y certificado en PDF.
\item La calificación usa el mismo semáforo de vencimiento a 30 y 15 días que los cursos.
\item Un soldador con la calificación vencida aparece bloqueado en el selector de personal del corte, con el motivo visible.
\item Las calificaciones no se editan ni se borran: se corrigen registrando una nueva.}

\historia{HU-R03}{Registrar los lotes de materiales que llegan al proyecto}{COULD}{3}
{Captura de cortes de obra.}
{Ingeniero Residente.}
{registrar los lotes de material que llegan al proyecto (por ejemplo, 4 cajas de soldadura) en lugar de lo que se gasta cada día.}
{saber qué material hay en la obra sin tener que calcular el consumo diario.}
{\item Puedo registrar la llegada de un lote con el material, la cantidad, la unidad y la fecha.
\item Cada lote queda asociado a su proyecto y a quien lo registró.
\item Puedo consultar los lotes recibidos por proyecto y por material.
\item El corte deja de pedir el consumo diario de los materiales que se controlan por lote.}

\historia{HU-R04}{Controlar herramientas, equipos y su mantenimiento}{SHOULD}{5}
{Equipos y herramientas.}
{Ingeniero Residente o Director de Proyectos.}
{anotar qué equipos y herramientas se usaron en el día y llevar su mantenimiento y sus certificados de calibración en un solo lugar, como se hace con las personas.}
{que ningún equipo trabaje con el mantenimiento o la calibración vencidos, aunque viaje entre obras.}
{\item En el corte puedo indicar qué equipos y herramientas se usaron en el día.
\item Cada equipo muestra su último mantenimiento y las horas que faltan para el siguiente.
\item El sistema avisa cuando falta poco para el siguiente mantenimiento y bloquea el equipo cuando ya venció, como el semáforo SST.
\item Cada equipo guarda sus certificados de calibración en PDF, con su vigencia.
\item Puedo ver en qué obra está cada equipo, aunque haya pasado por varias.}

\historia{HU-R05}{Hacer la revisión diaria de equipos y herramientas}{COULD}{3}
{Equipos y herramientas.}
{Director de Proyectos.}
{tener una lista diaria para comprobar que los equipos y herramientas de la obra están en condiciones de trabajar.}
{detectar a tiempo un equipo dañado antes de que afecte el trabajo o la seguridad.}
{\item Cada día hay una lista de revisión con los equipos asignados a la obra.
\item Por cada equipo marco si está apto o no apto, con una observación cuando no lo está.
\item Un equipo marcado como no apto no se puede seleccionar en el corte de ese día.
\item Queda registro de quién hizo la revisión y cuándo.}

\historia{HU-R06}{Anotar el clima del día en el corte}{COULD}{1}
{Captura de cortes de obra.}
{Ingeniero Residente.}
{anotar el clima del día en el corte (sol, lluvia y horas de lluvia).}
{explicar con datos por qué el avance fue menor un día.}
{\item El corte tiene un campo para el clima: soleado, nublado o lluvioso.
\item Si llovió, puedo indicar las horas de lluvia.
\item El clima se ve en el detalle del corte y en su historial de versiones.}

\historia{HU-R07}{Descargar el corte mensual para el cliente}{SHOULD}{3}
{Captura de cortes de obra.}
{Director de Proyectos o integrante de Gerencia.}
{recibir el corte mensual y descargarlo en un archivo para enviarlo al cliente.}
{entregar al cliente el avance del mes sin tener que armar el informe a mano.}
{\item Puedo armar el corte mensual a partir de los cortes visados del mes.
\item El corte mensual le llega al director o al gerente del proyecto.
\item Puedo descargarlo en PDF con el avance, las horas, los insumos y los soportes.
\item Solo se incluyen cortes visados.}

\historia{HU-R08}{Ver solo los cortes y documentos de mi obra}{SHOULD}{3}
{Control de acceso por roles.}
{Ingeniero Residente.}
{ver solo los cortes y documentos de la obra o ciudad en la que trabajo.}
{encontrar más rápido lo mío y que nadie consulte información de obras que no le corresponden.}
{\item Cada residente queda asignado a una o varias obras.
\item Las listas de cortes y documentos solo muestran los de sus obras.
\item Si abro por dirección un corte de otra obra, el sistema me avisa que no tengo permiso.
\item El director y la gerencia siguen viendo todas las obras.}

\historia{HU-R09}{Registrar un proyecto nuevo y avisar al director}{SHOULD}{2}
{Proyectos.}
{Ingeniero Residente.}
{registrar un contrato apenas se apruebe y que el director reciba un aviso.}
{empezar a reportar cortes en la obra sin esperar a que alguien lo cree por otro medio.}
{\item Un botón permite registrar un proyecto con su código, nombre, cliente, ciudad y fechas del contrato.
\item Al registrarlo, el director recibe un aviso.
\item El proyecto queda disponible para registrar cortes cuando el director lo confirma.}

\historia{HU-R10}{Recibir avisos por correo y recordatorios de facturación}{SHOULD}{3}
{Notificaciones.}
{Ingeniero Residente o Coordinación SIG / Calidad.}
{recibir un correo cuando me devuelvan un corte o haya que revisar el dossier, y un recordatorio cuando se acerque la facturación.}
{no depender de entrar al sistema para enterarme de lo que tengo pendiente.}
{\item Cuando el director devuelve un corte, el residente recibe un correo con la observación.
\item Cuando hay soportes por revisar en el dossier, Calidad recibe un correo.
\item Antes de la fecha de facturación, el residente recibe un recordatorio de los cortes pendientes.
\item Cada correo tiene un enlace directo a la pantalla correspondiente.}

\historia{HU-R11}{Que el encargado SST de cada obra actualice sus cursos}{COULD}{2}
{Control de habilitación SST.}
{Encargado SST de una obra.}
{actualizar yo mismo los cursos y certificados del personal de mi obra.}
{que la información SST esté al día sin depender de la coordinación central.}
{\item El encargado SST de una obra puede registrar cursos y certificados solo del personal de esa obra.
\item Los registros quedan con su autor, como los de la coordinación SST.
\item La coordinación SST sigue viendo y corrigiendo los registros de todas las obras.}

\end{document}
